\section*{Anexo G. Diagnóstico posterior del entrenamiento}
\phantomsection
\addcontentsline{toc}{subsection}{Anexo G. Diagnóstico posterior del entrenamiento}
\setcounter{table}{0}
\renewcommand{\thetable}{G.\arabic{table}}
\renewcommand{\theHtable}{G.\arabic{table}}
\setcounter{figure}{0}
\renewcommand{\thefigure}{G.\arabic{figure}}
\renewcommand{\theHfigure}{G.\arabic{figure}}
\setcounter{equation}{0}
\renewcommand{\theequation}{G.\arabic{equation}}
\renewcommand{\theHequation}{G.\arabic{equation}}
% Entrega 3, fase 2 (23.10.2026; plan del 14.10, movimientos 3, 5 y 6): este anexo junta el diagnostico posterior
% del entrenamiento que estaba en capitulo5.tex (l. 1401-1642 del commit 7181ab8, seccion "Analisis exploratorio:
% correccion del entrenamiento", cuya etiqueta sec:exploratorio se borro), capitulo3.tex (3.7.8,
% sec:diseno-exploratorio) y capitulo4.tex (4.7.7, sec:correccion-entrenamiento). El texto del Cap. V se movio sin
% reescribirlo: solo cambian las remisiones que ya no valen fuera del capitulo y se quitan los rotulos en negrita que
% repetian el titulo de su apartado. Pasada del glosario (23.10.2026, tarea 3f): IPPO con mensajes, ciclo minimo,
% fijo ajustado, version evaluada, diagnostico posterior del entrenamiento y fijo ajustado desde 10 s. Verificacion
% de la fase 2 (23.10.2026): tambien "cribado", "confirmacion", "examen", "brazo", "base" y "checkpoint" pasan a los
% nombres del glosario del plan del 14.10 (primera prueba, repeticion con semillas nuevas, evaluacion del plan,
% control o variante, entrenamiento y version guardada), y "Las dos correcciones" ya no explica dos veces seguidas
% el descuento y la exploracion. El material de
% los Caps. III y IV se integro el 23.10.2026 sin reescribirlo: cada tramo lleva un comentario "Integracion" con su
% origen y con lo que se ajusto (remisiones, la opcion de linea de ordenes, las rutas remitidas al Anexo B y los
% parrafos que repetian lo ya dicho en este anexo). Ninguna etiqueta despues de los titulos con asterisco; tablas y
% ecuaciones llevan su etiqueta despues del caption y numeran G.n.

Este anexo reúne el diagnóstico posterior del entrenamiento de IPPO con
mensajes: qué aprende y por qué, el diseño y su pre-registro, las dos
correcciones del entrenamiento, los resultados en P4, P6 y P1, lo que no
ayudó y sus límites. Su resumen, una frase sobre el origen de lo que
pasa en P4 y en P6, está en la Sección~\ref{sec:esperado-obtenido}. Se
remite también a este anexo desde el resumen de su diseño, al final de
la Sección~\ref{sec:diseno-robustez}, desde el de las correcciones en el
Capítulo~\ref{cap:implementacion}, desde lo que el plan no permite
afirmar, en el Capítulo~\ref{cap:resultados}, y desde las conclusiones
y las recomendaciones.

Este anexo es posterior al veredicto pre-registrado y
no lo modifica. Se hizo después de cerrar P6, y usa semillas de
evaluación distintas de las de la evaluación del plan: 5001 a 5030 para
una primera prueba y 6001 a 6030 para su repetición con semillas
nuevas. Las semillas 1001 a 1030 no se simularon. No debe confundirse con la fase exploratoria del
Anexo~C, que es anterior al plan. La pregunta es por qué IPPO con
mensajes no mejora la coordinación en P4 y P6, y si eso se corrige cambiando
solo su entrenamiento. El diseño está más abajo, en ``Diseño y
pre-registro'', y los cambios del entrenamiento, en ``Las dos
correcciones''. La métrica es la del Capítulo~\ref{cap:resultados}, el
\texttt{timeLoss} medio por vehículo, y la estadística, la de la
Sección~\ref{sec:estadistica}.
% fuente: plan2/exploratorio_coord/LEEME.md y HECHOS.md (estatus; cada script lleva un assert contra 1001-1030);
% plan2/exploratorio_coord/PRERREGISTRO.md (alcance y semillas, seccion 7)

% Integracion (23.10.2026): aqui iba capitulo3.tex l. 1611-1614 del commit 7181ab8 (comienzo de 3.7.8,
% sec:diseno-exploratorio). No se repite: lo que dice (despues del analisis de robustez se estudio por que IPPO
% vecinos no mejora la coordinacion en P4 y P6; es posterior al veredicto, usa otras semillas y no cambia ningun
% veredicto) ya esta en los dos parrafos de arriba. El texto sigue en el commit 7181ab8.

\subsection*{Qué aprende IPPO con mensajes y por qué}

\noindent\textbf{Qué aprende hoy IPPO con mensajes.} Su versión
evaluada no aprende lo mismo en todos los casos. En P3 y P5 aprende el ciclo
mínimo y lo sincroniza con los mensajes: si se cortan todos los
mensajes, su \texttt{timeLoss} sube un $79.8$\,\% en P3 y un
$70.1$\,\% en P5 (semillas 999 y 1999). En P4 aprende un plan casi
fijo y largo, de unos 25~s de avenida y 10~s de transversal, con un
ciclo de 41~s; ahí cortar los mensajes no lo empeora ($-4.1$\,\%). En
P6 cortarlos sí lo empeora ($+131.8$\,\%), pero el programa que
aprende también es largo y poco sincronizado. Con la semilla 999, su
verde de avenida tiene una mediana de 25~s, su ciclo una de 46~s, y
dos cruces vecinos están en la misma fase solo el $18.3$\,\% de la
hora.
% fuente: plan2/exploratorio_coord/diagnostico/mejora_coord/a3_obs_sensibilidad.txt (perdida 1.0: red +79.8,
% malla3 +70.1, red_alta -4.1, malla3_alta +131.8); diagnostico/mejora_coord/SINTESIS.md (P4 ~24-25/10, ciclo
% 41 s); sintesis/analizar_fases_p6.txt (ippo_s42: vEO 25, ciclo 46, sinc 18.3 %)

\noindent\textbf{Por qué.} El diagnóstico apunta a dos causas en el
entrenamiento. La primera es el descuento. GAE
(Ecuación~\ref{eq:gae}) descuenta por decisión, y las decisiones no
duran lo mismo: 5~s si se mantiene la fase y 8~s si se cambia, ámbar
incluido (Sección~\ref{sec:aprendices}). Medido en decisiones, un
verde largo parece más barato de lo que es en segundos. En P4, con
ocho semillas de validación, la versión evaluada tiene un
\texttt{timeLoss} un $4.8$\,\% mayor que el del ciclo mínimo, pero su
costo descontado por decisión es un $1.5$\,\% menor: el objetivo del
entrenamiento prefiere el plan largo aunque el tráfico espere más. En
P6 pasa lo mismo frente al fijo ajustado: $+4.6$\,\% en
\texttt{timeLoss} y $-1.6$\,\% en el costo por decisión. La
segunda es la exploración. Cada semáforo sortea su acción por
separado, y un cambio de fase a destiempo en un solo cruce rompe el
ciclo común. Con un 3\,\% de acciones al azar sorteadas por separado,
el ciclo mínimo de P4 empeora un $21.7$\,\% y el fijo de P6, un
$14.7$\,\%. Si todos los semáforos usan el mismo número aleatorio, el
daño baja a $1.2$ y $2.1$\,\%. Durante el entrenamiento, ese ruido
castiga a los programas que dependen de la sincronía.
% fuente: plan2/exploratorio_coord/diagnostico/c4_alineacion.txt (P4, val8, ippo_s42 frente al trivial: timeLoss
% +4.8 %, J_dec -1.5 %) y c4_alineacion_p6.txt (P6 frente al 15/10: +4.6 %, J_red_dec -1.6 %);
% diagnostico/b2_ruido.txt (val8: red_alta trivial 27.08 -> 32.96 independiente, 27.39 comun; malla3_alta 15/10
% 36.57 -> 41.94 y 37.36); duracion de las decisiones: PRERREGISTRO.md (V2) y docstring de gae()

% Integracion (23.10.2026): aqui iba capitulo3.tex l. 1615-1624 del commit 7181ab8 (3.7.8: las dos causas que
% senalo un diagnostico previo, sin entrenar). No se repite: las dos causas, el descuento por decision con
% decisiones de 5 y 8 s (eq:gae, sec:aprendices) y la exploracion sorteada por separado en cada semaforo, ya estan en
% el parrafo anterior con sus cifras, y su remision a sec:exploratorio remitia a este mismo apartado. El texto sigue
% en el commit 7181ab8.

\subsection*{Diseño y pre-registro}

% Integracion (23.10.2026): capitulo3.tex l. 1626-1713 del commit 7181ab8 (3.7.8), con la Tabla
% tab:exploratorio-variantes (aqui G.1) y sus fuentes. Sin reescribir; solo cambian las remisiones: la de l. 1636 a
% sec:correccion-entrenamiento pasa a "mas abajo, en Las dos correcciones"; la de la fila V2 de la tabla (l. 1650) se
% quita, y la de l. 1708 a sec:exploratorio pasa a "mas abajo, en Resultados". En la pasada del glosario
% (23.10.2026) el texto y el caption pasaron a los nombres nuevos. El parrafo que resumia este mismo diseno desde capitulo5.tex (l.
% 1474-1490 del commit 7181ab8) no se repite: cada una de sus frases esta en el texto siguiente (C0 y las cinco
% variantes, tres bases y 200 episodios, la regla de seleccion y convergencia del plan, ganadora en 5001-5030 y
% confirmacion sin reentrenar en 6001-6030, las dos mejores a P6 y la ganadora a P1, la confirmacion de P6 fuera del
% pre-registro, P2, P3 y P5 sin probar). Su fuente (PRERREGISTRO.md secciones 2 a 7; prerregistro.sha256
% 8d38fe82...6b33; DESVIACIONES.md D1-D11 y T1-T7, replica de C0 en D4 y de C0-P6 en T7; conf_p6/juzgar_conf_p6.py,
% criterio CONFIRMA; HECHOS.md seccion 1) es la misma de los parrafos siguientes.
El diseño se fijó en un pre-registro propio antes de entrenar,
con su huella sha256 guardada aparte. Los cambios
posteriores se declararon como desviaciones, sin editar el
pre-registro. Se comparan un control y cinco variantes
(Tabla~\ref{tab:exploratorio-variantes}). El control, C0, es IPPO con
mensajes sin cambios en el código, con la validación de ocho semillas
que usan las variantes. C0, V1, V2 y V3 forman un diseño factorial de
dos factores: la observación normalizada y las dos correcciones del
entrenamiento. V4 y V5 agregan a V3 un solo factor cada una. La
implementación de las dos correcciones está más abajo, en ``Las dos
correcciones''.

\begin{table}[H]
\setstretch{1}
\centering
\small
\begin{tabular}{l>{\raggedright\arraybackslash}p{10.8cm}}
\toprule
Control o variante & Cambio respecto de IPPO con mensajes del plan \\
\midrule
C0 & Ninguno en el código; validación con ocho semillas \\
V1 & Observación normalizada: colas y vehículos en movimiento de cada
aproximación divididos entre 10~veh por carril, con tope 2 \\
V2 & Descuento por segundo en la GAE y exploración con un número
aleatorio común \\
V3 & V1 y V2 juntas \\
V4 & V3 con MAPPO y una recompensa de toda la red, que incluye los
enlaces sin semáforo \\
V5 & V3 con el actor iniciado en el ciclo mínimo: probabilidad
inicial de cambiar cercana a 0.97 en toda decisión libre \\
\bottomrule
\end{tabular}
\caption[Diagnóstico posterior: control y variantes]{Diagnóstico posterior del entrenamiento:
control y variantes de IPPO con mensajes. Todos usan los hiperparámetros de la Tabla~\ref{tab:plan2-hiper}, tres
entrenamientos y 200 episodios.}
\label{tab:exploratorio-variantes}
\end{table}
% fuente: plan2/exploratorio_coord/PRERREGISTRO.md (secciones 2 y 3; sha256 en prerregistro.sha256:
% 8d38fe82...6b33) y DESVIACIONES.md (D1-D11, T1-T7); HECHOS.md seccion 1

El resto del protocolo es el del plan
(Sección~\ref{sec:protocolo}). El control y cada variante se entrenan
tres veces, con las semillas 42, 2042 y 4042 y 200 episodios, con los
hiperparámetros de la Tabla~\ref{tab:plan2-hiper}. La elección de la
versión guardada de cada entrenamiento, la de la versión evaluada y el
criterio de convergencia son los mismos, y el episodio 0 no es
candidato. Solo cambia la validación, que usa ocho
semillas en lugar de dos: 999, 1999, 2999, 3999, 4999, 7999, 8999 y
9999. Se valida cada 10 episodios en P4 y en P1 y cada 20 en P6. C0 no
se reentrena. Como la validación guarda y restaura el estado del
generador aleatorio, entrenar de nuevo daría los mismos pesos, así que
C0 se obtiene validando con las ocho semillas las versiones ya
guardadas del plan. Antes se comprobó que diez episodios
reentrenados reproducen los pesos guardados.

La evaluación usa la política congelada, el \texttt{timeLoss} y la
estadística pareada de la Sección~\ref{sec:estadistica}. La regla de
la Sección~\ref{sec:veredicto} se aplica frente al fijo del resultado
principal. Ninguna evaluación usa las semillas de la evaluación del
plan, 1001 a 1030, y cada script lo comprueba. La primera prueba usa las
semillas 5001 a 5030 y la repetición, 6001 a 6030, con las mismas
políticas y sin reentrenar. Las semillas 6001 a 6010 ya se habían simulado en el
diagnóstico, así que se informa además el resultado en 6011 a 6030.

El análisis tiene dos tandas. La primera, en P4, entrena las cinco
variantes. Gana la de menor \texttt{timeLoss} en 5001 a 5030 entre las
que convergen. Para declarar un éxito, la ganadora debía mejorar al
fijo en al menos un 5\,\%, con intervalo entero bajo cero y $p < 0.05$,
superar al ciclo mínimo con intervalo entero bajo cero y repetir
las dos cosas en 6001 a 6030. La segunda tanda lleva las dos mejores
variantes de P4, sin cambios, a P6, junto con el control C0 de ese
caso. En P6 una variante mejora si el entrenamiento de su versión
evaluada converge, si el intervalo de su diferencia frente a IPPO con
mensajes del plan queda entero bajo cero y si la regla frente al fijo
no da FALLA. En P1 la ganadora debe dar FUNCIONA, y el límite superior
de su intervalo frente a IPPO con mensajes del plan no debe pasar de
$+2$\,\%. La repetición de P6 en 6001 a 6030 no estaba en el
pre-registro. Se agregó después de ver la primera prueba de P6, con el
criterio de que las dos versiones evaluadas mejoren al fijo y a IPPO
con mensajes del plan con intervalo entero bajo cero. P2, P3 y P5 no se prueban.

La ganadora se elige en 5001 a 5030, así que sus cifras en esas
semillas son optimistas; solo las de 6001 a 6030 confirman. Los
resultados están más abajo, en ``Resultados''.
% fuente: plan2/exploratorio_coord/PRERREGISTRO.md (seccion 2: bases, val8, eleccion, convergencia; seccion 4:
% ganadora y EXITO a, b, c; seccion 6: tanda 2, criterio de mejora en P6, NO REGRESION en P1 con +2 %,
% sensibilidad 6011-6030; seccion 7: semillas 1001-1030 prohibidas); DESVIACIONES.md D4 y T7 (replica de C0 y
% C0-P6 por re-validacion); conf_p6/juzgar_conf_p6.py (docstring: criterio CONFIRMA, fuera del pre-registro);
% LEEME.md (assert contra 1001-1030)

\subsection*{Las dos correcciones}

% Verificacion de la fase 2 (23.10.2026): los dos parrafos de entrada de este apartado (el que venia del Cap. V,
% capitulo5.tex l. 1460-1472 del commit 7181ab8, y el de la 4.7.7) explicaban dos veces seguidas el descuento y la
% exploracion antes del detalle que sigue; quedan en uno, que remite a las causas de "Que aprende IPPO con mensajes
% y por que" (antes remitia a "Diseno y pre-registro", donde no estan).
% Integracion (23.10.2026): capitulo4.tex l. 1600-1680 del commit 7181ab8 (4.7.7), con la Ecuacion eq:gae-smdp (aqui
% G.1) y sus fuentes. Sin reescribir; solo cambian: (1) la remision a sec:diseno-exploratorio (l. 1602) pasa a "mas
% arriba, en Diseno y pre-registro"; (2) sale la opcion de linea de ordenes --variante smdp_crn (l. 1605); (3) las
% rutas de plan2/exploratorio_coord/ (l. 1606-1608 y 1666-1671: rl_ippo_smdp_crn.py, smdp_crn.diff,
% rl_ippo_norm_smdp_crn_triv.py y lo que guarda la carpeta, con verificar_hechos.py y la remision a sec:exploratorio)
% se remiten al Anexo B, que las lista; (4) no se repite el parrafo de l. 1655-1661 (PPO con el logaritmo de cada
% accion por separado, estimador sesgado, heuristica, los dos cambios siempre juntos), porque el sesgo y la
% heuristica estan en el parrafo siguiente y los dos cambios juntos en "Lo que no ayudo".
\begin{sloppypar}
La variante V2 cambia los dos puntos del entrenamiento de IPPO que
señala el diagnóstico (más arriba, en ``Qué aprende IPPO con mensajes y
por qué''), el descuento de la ventaja y la forma de explorar, y nada
más. Su código y su diferencia con \texttt{rl\_ippo.py}, el único
archivo que cambia, están en la carpeta del análisis (Anexo~B). La
observación, la máscara, la recompensa, los mensajes y los
hiperparámetros son los del plan. La evaluación tampoco cambia: con la
política congelada se toma la acción de mayor probabilidad, y ninguna
de las dos correcciones actúa fuera del entrenamiento. La exploración
común es una heurística: con probabilidades distintas entre semáforos,
el estimador del gradiente queda sesgado.
\end{sloppypar}
% fuente: plan2/exploratorio_coord/codigo/smdp_crn.diff y rl_ippo_smdp_crn.py (gae con gamma^(dt/5); u comun de
% default_rng(semilla + 104729)); limitacion del estimador: PRERREGISTRO.md (V2) y SINTESIS.md seccion 6;
% HECHOS.md seccion 1 (codigo cambiado: solo rl_ippo.py)

La primera corrección es el descuento por segundo. En el plan, un
intervalo de 5~s y uno de 8~s cuentan como un paso con el mismo
descuento (Sección~\ref{sec:aprendices}). En la corrección cada
decisión guarda el segundo $t_k$ en que se tomó, y la duración de su
intervalo es $\Delta t_k = t_{k+1} - t_k$; en la última decisión,
$t_{k+1}$ es el segundo en que cierra el episodio. La
Ecuación~\ref{eq:gae} se calcula entonces hacia atrás con exponentes
proporcionales a esa duración:
\begin{equation}
  \delta_k = r_k + \gamma^{\Delta t_k/5}\, V(o_{k+1}) - V(o_k), \qquad
  \hat A_k = \delta_k + (\gamma\lambda)^{\Delta t_k/5}\, \hat A_{k+1},
  \label{eq:gae-smdp}
\end{equation}
donde $r_k$ es la recompensa de la decisión $k$, que no cambia, y
$V(o_k)$ el valor que estima el crítico. Si se mantiene la fase el
exponente es 1, como en el plan; si se cambia es 1.6. El valor final,
cero si la red se vació y el estimado si el episodio se cortó, se
descuenta igual.

La segunda corrección es la exploración común. En el plan cada
semáforo muestrea su acción de su propia distribución, de forma
independiente. En la corrección se sortea un solo número $u$, uniforme
entre 0 y 1, en cada segundo en que algún semáforo decide, y es el
mismo para todos los que deciden en ese segundo. Cada uno cambia de
fase si su probabilidad de cambiar, calculada con la política
enmascarada, supera a $u$. La probabilidad de cambiar de cada
semáforo sigue siendo la de su política, pero sus desvíos quedan
correlacionados: con probabilidades parecidas, los semáforos se
desvían juntos y el ciclo común se conserva. La máscara sigue
decidiendo sola cuando una de las dos acciones está prohibida. El
número $u$ sale de un generador propio, \texttt{default\_rng} con la
semilla del episodio más 104\,729, que no toca los
generadores de Python, NumPy ni PyTorch; el resto del entrenamiento y
la validación usan los mismos números que en el plan.
% fuente: plan2/exploratorio_coord/codigo/smdp_crn.diff (correr_episodio: crn, rng_expl, u, p_cambiar, acciones,
% logp, b["t"]; intervalos; gae con smdp=True; SEMILLA_EXPL = 104729) y rl_ippo_smdp_crn.py (docstring de gae:
% 5 s mantener, 8 s cambiar, t_final - t en la ultima, tambien para v_final); PASO_CONTROL = 5 (rl_corredor.py)
% 22.10.2026 (entrega 3, sin codigo en la tesis): se quito el extracto simplificado de correr_episodio y gae con los
% dos cambios; el texto y la Ecuacion eq:gae-smdp los describen. Sigue en el commit 5b9b7a6

V5 agrega a V2 la observación normalizada de V1 y un actor que empieza
en el ciclo mínimo. Su código, necesario también para cargar sus
políticas, está en la carpeta del análisis (Anexo~B). Los metadatos de cada política declaran el
descuento y la exploración con que se entrenó. Los entrenamientos
usan SUMO 1.25.0 y un hilo por proceso, y son deterministas: la
semilla 5001, repetida en la repetición con semillas nuevas, dio lo
mismo en el control y en todas las variantes.
% fuente: plan2/exploratorio_coord/LEEME.md (contenido, como se reproduce, versiones), HECHOS.md seccion 1 (codigo
% cambiado: solo rl_ippo.py; sha256 iguales a tanda2/lanzamiento_p6_entrenamiento.sha256) y seccion 8 (limitaciones
% 4 y 5); HECHOS.md secciones 4 y 5.2 (determinismo, semilla 5001 en los 6 y 10 brazos); smdp_crn.diff (meta
% "descuento" y "exploracion"); politicas/*.json

\subsection*{Resultados}

% fuente: plan2/exploratorio_coord/tanda1/juicio_tanda1.txt (P4, 5001-5030), tanda2/juicio_tanda2.txt (P4
% 6001-6030, P6 5001-5030 y P1), conf_p6/salidas/juicio_conf_p6.txt (P6 6001-6030 y 60 semillas); todas las
% cifras recalculadas desde los CSV por semilla con plan2/exploratorio_coord/verificar_hechos.py (verificar_hechos.txt)
% Integracion (23.10.2026): flota ([!tbp], tabla grande, regla de formato del plan); en [H] dejaba en blanco media
% pagina antes de "Resultados"
\begin{table}[!tbp]
\setstretch{1}
\centering
\footnotesize
\setlength{\tabcolsep}{3pt}
\begin{tabular}{lrrrr}
\toprule
\makecell[l]{Control o variante\\(entrenamiento)} & \makecell{Retraso\\(s/veh)} & \makecell{Frente al fijo\\ajustado (\%)} & \makecell{Frente al\\ciclo mínimo (\%)} & \makecell{Frente a IPPO\\del plan (\%)} \\
\midrule
\multicolumn{5}{l}{\textit{P4, primera prueba, semillas 5001 a 5030}} \\
C0 (4042) & 29.08 & $-3.3$ [$-4.8$, $-1.8$] & $+3.3$ [$+1.6$, $+5.0$] & $+0.5$ [$-0.9$, $+2.0$] \\
V2 (2042) & 28.15 & $-6.4$ [$-8.4$, $-4.3$] & $0.0$ (igual) & $-2.7$ [$-4.4$, $-1.0$] \\
V5 (42) & 28.07 & $-6.6$ [$-8.6$, $-4.7$] & $-0.3$ [$-0.8$, $+0.2$] & $-3.0$ [$-4.7$, $-1.3$] \\
\multicolumn{5}{l}{\textit{P4, repetición con semillas nuevas, 6001 a 6030}} \\
IPPO del plan (42) & 28.74 & $-3.1$ [$-4.8$, $-1.5$] & $+3.3$ [$+1.5$, $+5.1$] & -- \\
V5 (42) & 27.83 & $-6.2$ [$-7.8$, $-4.5$] & $+0.0$ [$-0.3$, $+0.4$] & $-3.2$ [$-4.9$, $-1.5$] \\
\multicolumn{5}{l}{\textit{P6, primera prueba, semillas 5001 a 5030}} \\
IPPO del plan (42) & 37.75 & $+2.4$ [$-1.7$, $+7.1$] & -- & -- \\
C0 (4042) & 39.64 & $+7.6$ [$+2.0$, $+13.3$] & -- & $+5.0$ [$-1.4$, $+11.8$] \\
V2 (4042) & 34.14 & $-7.3$ [$-10.0$, $-5.3$] & -- & $-9.5$ [$-13.8$, $-5.4$] \\
V5 (42) & 34.22 & $-7.1$ [$-9.8$, $-4.8$] & -- & $-9.3$ [$-13.8$, $-5.1$] \\
\multicolumn{5}{l}{\textit{P6, repetición con semillas nuevas, 6001 a 6030}} \\
IPPO del plan (42) & 37.59 & $+1.3$ [$-2.7$, $+5.5$] & -- & -- \\
V2 (4042) & 34.65 & $-6.7$ [$-9.8$, $-3.7$] & -- & $-7.8$ [$-11.7$, $-3.8$] \\
V5 (42) & 34.55 & $-6.9$ [$-10.3$, $-3.7$] & -- & $-8.1$ [$-11.8$, $-4.4$] \\
\multicolumn{5}{l}{\textit{P6, las 60 semillas (5001 a 5030 y 6001 a 6030)}} \\
IPPO del plan (42) & 37.67 & $+1.9$ [$-1.1$, $+5.0$] & -- & -- \\
V2 (4042) & 34.40 & $-7.0$ [$-9.0$, $-5.1$] & -- & $-8.7$ [$-11.7$, $-5.7$] \\
V5 (42) & 34.38 & $-7.0$ [$-9.2$, $-5.0$] & -- & $-8.7$ [$-11.6$, $-5.8$] \\
\multicolumn{5}{l}{\textit{P1, no regresión, semillas 5001 a 5030}} \\
IPPO del plan (4042) & 19.85 & $-6.7$ [$-8.0$, $-5.4$] & -- & -- \\
V5 (4042) & 19.73 & $-7.3$ [$-8.7$, $-6.0$] & -- & $-0.6$ [$-1.5$, $+0.3$] \\
\bottomrule
\end{tabular}
\caption[Diagnóstico posterior: variantes del entrenamiento de IPPO con mensajes]{Diagnóstico posterior del entrenamiento, fuera del veredicto. Retraso (\texttt{timeLoss} medio, s/veh) de la versión evaluada del control y de cada variante y su cambio porcentual frente al fijo ajustado del resultado principal (en P6, 15~s de avenida y 10~s de transversal sin desfases), al ciclo mínimo y a IPPO con mensajes del plan, con el IC95\,\% por \textit{bootstrap} pareado entre corchetes. C0 es IPPO con mensajes sin cambios validado con ocho semillas; V2, descuento por segundo y exploración común; V5, V2 con la observación normalizada y el inicio en el ciclo mínimo. IPPO del plan es la versión evaluada de IPPO con mensajes en el plan. ``Igual'': el mismo resultado que el ciclo mínimo en las 30 semillas. En P1 y P6 no se da el cambio frente al ciclo mínimo, que falla por mucho (102.58 y 184.39~s/veh en 5001 a 5030). Medias de las referencias en P4: fijo 30.06 y ciclo mínimo 28.15 (5001 a 5030); fijo 29.66 y ciclo mínimo 27.82 (6001 a 6030). En P6 el fijo da 36.84 y 37.12, y en P1, 21.27. Un valor negativo es una mejora.}
\label{tab:exploratorio}
\end{table}

La Tabla~\ref{tab:exploratorio} resume los resultados por caso. Los
archivos por semilla y el script que recalcula sus cifras están en la
carpeta del análisis (Anexo~B).

\noindent\textbf{P4: el ciclo mínimo.} Con el descuento por segundo y
la exploración común, IPPO con mensajes aprende en P4 el ciclo mínimo.
La versión evaluada de V2 da el mismo resultado que el ciclo mínimo
(Sección~\ref{sec:programa-trivial}) en las 30 semillas, y la de V4
también. V5 fue la ganadora de la primera prueba: $-6.6$\,\% frente al
fijo ajustado y $-0.3$\,\% frente al ciclo mínimo, con un intervalo que
incluye el cero. En la repetición con semillas nuevas queda en
$-6.2$\,\% y $+0.0$\,\%, así que la pequeña ventaja de la primera
prueba venía, probablemente, de elegir a la
mejor de cinco variantes con esas mismas semillas. El pre-registro pedía, para declarar un éxito, ganarle
también al ciclo mínimo, y eso no se cumple: V5 reproduce el ciclo
mínimo y no lo supera. Frente a IPPO con mensajes del plan reduce el
retraso un $3.2$\,\% en la repetición. Frente al fijo ajustado desde
10~s del Anexo~F, que en P4 es un ciclo de 10~s por fase con
desfases, V5 empata: $+1.6$\,\%, de $-0.2$ a $+3.5$\,\%, en las
semillas 5001 a 5030; ese fijo no se simuló en 6001 a 6030. La
ganancia en P4 consiste en pasar del plan largo al ciclo corto, algo
que un fijo bien barrido también logra.
% fuente: tanda1/juicio_tanda1.txt (V2 y V4 igual_trivial 30/30; V5 ganadora; criterio (b) no se cumple);
% tanda2/juicio_tanda2.txt seccion 1 (V5 6001-6030; NO hay EXITO); sintesis/extra_robustez.txt (fijo re-barrido
% 10/10, desfases 0/25/24: V5 +1.6 % [-0.2, 3.5], p 0.118, gana 13/30)

\noindent\textbf{P6: una mejora real.} En P6 el resultado es otro. La
versión evaluada de V2 reduce el retraso un $7.3$\,\% frente al fijo
ajustado en la primera prueba y gana en 28 de las 30 semillas. En la
repetición con semillas nuevas lo reduce un $6.7$\,\%, de $-9.8$ a $-3.7$\,\%, y gana en
29 de 30. Frente a IPPO con mensajes del plan, que en esas mismas semillas
empata con el fijo ($+1.3$\,\%), como en la evaluación del plan ($+2.2$\,\%), la
reducción es de un $7.8$\,\%. V5 da lo mismo: $-7.1$ y $-6.9$\,\%. El
pre-registro no esperaba ganarle al fijo en P6, por eso la mejora se
revisó con tres cortes más. Con las 60 semillas, V2 queda en
$-7.0$\,\% (de $-9.0$ a $-5.1$\,\%). En las semillas 6011 a 6030, que
no se habían usado antes en ningún diagnóstico, queda en $-6.5$\,\%
(de $-11.2$ a $-2.1$\,\%). Sin la semilla 6028, en la que el fijo da
50.48~s/veh, su peor valor en esas 30 semillas (el siguiente es
38.01~s/veh), queda en $-5.4$\,\% (de $-7.1$
a $-3.1$\,\%). Es la primera vez en esta tesis que un control le gana
de forma consistente al fijo ajustado de P6.
% fuente: tanda2/juicio_tanda2.txt seccion 2 (V2 4042: -7.3 % [-10.0, -5.3], 28/30); conf_p6/salidas/
% juicio_conf_p6.txt (6001-6030: V2 29/30, IPPO publicado +1.3 %; 60 semillas); verificar_hechos.txt (6011-6030 y
% sin 6028; siguiente valor del fijo 38.01); examen: plan2/veredicto.csv (malla3_alta, ippo_vecinos +2.2)

Lo que aprende en P6 no es el ciclo mínimo, que ahí se traba
(184.39~s/veh), ni el fijo de 15 y 10~s. Con la semilla 999, V2 y V5
parten del ciclo mínimo: su ciclo tiene una mediana de 26~s, la
transversal recibe siempre 10~s, y el $72$\,\% (V2) y el $66$\,\% (V5)
de los verdes de avenida duran 10~s. El resto se extiende a entre 15 y
35~s, con verdes distintos dentro de un mismo cruce. Dos cruces
vecinos están en la misma fase el $63.6$\,\% y el $60.5$\,\% de la
hora, frente al $18.3$\,\% de IPPO con mensajes del plan. No se sabe a qué
reacciona la extensión del verde, y no se comparó con una regla simple
de ciclo mínimo con extensión cuando hay cola, que podría igualarla.
Por eso el resultado se describe como una política adaptativa que
supera a todas las referencias medidas, y no como una coordinación
aprendida demostrada. En convergencia, los entrenamientos de las
versiones evaluadas de V2 y de V5 convergen, pero el de la semilla 2042
no lo hace en ninguna de las dos, y la versión evaluada de V5 se eligió
en el último episodio. C0 no converge en P6 en ninguno de sus tres
entrenamientos.
% fuente: sintesis/analizar_fases_p6.txt (semilla 999: ciclo 26, vNS 10, %vEO=10 72 y 66, verdes E-O de 15 a 35 s,
% sinc 63.6/60.5/18.3); SINTESIS.md seccion 3.1 (lo que no se sabe); tanda2/juicio_tanda2.txt (convergencia por
% base; V5 ep_elegido 200; C0-P6 0/3)

\noindent\textbf{P1: sin regresión.} En el corredor, V5 queda en
$-7.3$\,\% frente al fijo ajustado, e IPPO con mensajes del plan, en
$-6.7$\,\%. La diferencia entre los dos, $-0.6$\,\% con un límite
superior de $+0.3$\,\%, cumple el umbral de no regresión del
pre-registro, $+2$\,\%. La corrección no daña el caso en que IPPO con
mensajes ya funcionaba. V2 no se corrió en P1.
% fuente: tanda2/juicio_tanda2.txt seccion 3 (V5 4042 frente a ippo_vecinos_s4042 -0.6 % [-1.5, 0.3], p 0.289;
% NO REGRESION: SI)

\subsection*{Lo que no ayudó}

Lo que mejora el resultado es la corrección de V2: en P4, V2 reduce el
retraso un $3.2$\,\% frente a C0 (de $-4.8$ a $-1.6$\,\%, $p = 0.006$
con la corrección de Holm). Normalizar las colas, sola, no cambia nada:
V1 queda a $0.0$\,\% de C0 (de $-1.7$ a $+1.6$\,\%). MAPPO con la
recompensa de toda la red (V4) termina en el ciclo mínimo. Elegir la
versión guardada con ocho semillas en vez de dos tampoco ayuda: C0
queda a $+0.5$\,\% de IPPO con mensajes del plan en P4 y a $+5.0$\,\%
en P6, los dos con un intervalo que incluye el cero. En P6, el inicio
en el ciclo mínimo y la observación normalizada de V5 no suman nada a
V2 ($+0.2$\,\% en la primera prueba y $-0.3$\,\% en la repetición).
Como V2 cambia el descuento y la exploración a la vez, no se sabe cuál
de los dos pesa más.
% fuente: tanda1/juicio_tanda1.txt (contrastes frente a C0 con Holm; C0 frente a IPPO publicado); tanda2/
% juicio_tanda2.txt (C0-P6 +5.0 % [-1.4, 11.8]; V5 frente a V2 +0.2 %); conf_p6/salidas/juicio_conf_p6.txt (-0.3 %)

\subsection*{Límites}

El análisis es post hoc: se decidió
después de ver el veredicto, y la repetición de P6 con semillas nuevas
se agregó después de la primera prueba. Con el retraso total de la Ecuación~\ref{eq:retraso-total}
la mejora de P6 se reduce a menos de la mitad: V2 queda en $-3.1$\,\% (de
$-4.7$ a $-1.7$\,\%) frente al fijo en 6001 a 6030. La espera para
entrar a la red, de unos 45.5~s/veh, es la misma con todos los
controles salvo el ciclo mínimo, así que no es cola desplazada fuera de la
red, pero diluye la ganancia. En P4 esa espera es de 0.34~s/veh y el
retraso total da la misma lectura. Las conclusiones descansan en tres
entrenamientos por variante, y en P6 uno de ellos no converge. La demanda es
determinista, así que las semillas solo cambian el comportamiento de
los conductores. P2, P3 y P5 no se probaron, y en P4 la mejora
desaparece frente al fijo ajustado desde 10~s. Ningún veredicto del
Capítulo~\ref{cap:resultados} cambia. Lo que el análisis sí muestra es que el plan largo y
poco sincronizado que IPPO con mensajes aprende en P4 y P6 se debe en
buena parte a cómo se entrena, y que corregir el descuento y la
exploración basta para que en P6 supere al fijo ajustado. Medirlo con el protocolo
completo, en los seis casos, queda como trabajo siguiente.
% fuente: conf_p6/salidas/juicio_conf_p6.txt (retraso total V2 4042 -3.1 % [-4.7, -1.7]; retraso de insercion
% 45.49-45.65 s/veh salvo el trivial 95.51; 0 no insertados); sintesis/extra_robustez.txt (P4: insercion 0.336 s/veh,
% V5 retraso total -6.6 % [-8.5, -4.6]); tanda2/juicio_tanda2.txt (convergencia); HECHOS.md seccion 8
